\chapter{Literature Review}\label{chap:survey}

\section{Systematic Review - Deep Learning and Radiomics in Renal Tumour Classification}
\textbf{This systematic review is currently being prepared for submission to a peer-reviewed journal.}

\vspace{2mm}

Before presenting the findings of the systematic review, it is important to outline its position within the context of this research project. This review was conceived as an initial step to critically assess the current landscape of ML and DL applications in renal tumour texture analysis before answering the research questions. Tumour classification tasks, such as distinguishing between benign and malignant renal tumours, or subtyping RCC, serve as a relevant and clinically meaningful entry point for evaluating the potential of DL-based texture analysis.


The review aimed to determine how DL compares with traditional radiomics-based ML methods in terms of diagnostic performance, methodological consistency, and data requirements. Classification was chosen for its direct clinical relevance and reliance on texture information, making it an appropriate test case for assessing DL feasibility. The findings showed that radiomics combined with ML remains the dominant paradigm, while evidence for DL is present but limited. This confirmed that DL is feasible for texture analysis but highlighted a clear research gap: DL is not yet widely applied to renal tumour classification, and few studies directly compare DL with radiomics-based ML pipelines. The review also exposed substantial methodological heterogeneity across studies, including variations in CT acquisition phases, feature extraction procedures, model architectures, and evaluation strategies. Such inconsistencies emphasise the need for more rigorous and standardised methodologies. Notably, the absence of studies combining radiomics features with DL representations underscores the potential value of multimodal fusion, motivating a shift in research focus towards unified fusion frameworks.


\subsection{Results}\label{subsec:result}
            
\begin{figure}[ht]
 \centering
 \makebox[\textwidth][c]{\includegraphics[width=0.75\textwidth]{pictures/flow.png}}%
 \caption{PRISMA Flow-chart of included studies. There are 24 studies included.}
 \label{fig-1}
\end{figure}
            
A total of $274$ records were identified through database searches ($250$ English-language and $24$ Chinese-language articles). Chinese publications were included to ensure comprehensive coverage of rapidly expanding radiomics and DL research from China, with screening supported by a Chinese-speaking author. After removing duplicates, $180$ records remained; title and abstract screening excluded $100$, leaving $80$ for full-text review. Following application of all eligibility criteria, $24$ studies were included. These focused specifically on CT-based radiomics or DL methods for renal tumour classification and reported relevant performance metrics. The selection process is summarised in Figure~\ref{fig-1}.

\subsection{Private vs Public Datasets}

Most studies (79\%) relied exclusively on internal institutional datasets. While such datasets benefit from consistent imaging protocols and expert annotation, they limit generalisability and reproducibility. Public datasets are therefore essential for benchmarking and external validation. However, heterogeneity in annotation standards, imaging phases, and data quality persists even within public resources, underscoring the need for larger, well-curated repositories with consistent metadata.

\subsection{Public Data Source}\label{Sec:dataset}

Nineteen studies used internal datasets, whereas only five incorporated public datasets. The primary public resources were:
\begin{itemize}
    \item The Cancer Genome Atlas (TCGA),
    \item The Cancer Imaging Archive (TCIA)\footnote{\href{https://www.cancerimagingarchive.net/}{https://www.cancerimagingarchive.net/}}, and
    \item The Kidney Tumor Segmentation Challenge (KiTS) datasets.
\end{itemize}

TCGA provides multi-omics profiles and extensive clinical information across numerous tumour types~\cite{weinstein2013cancer}, but no imaging data. TCIA complements TCGA by offering de-identified CT, MRI, and PET imaging~\cite{clark2013cancer}, including collections that bridge imaging and genomics (e.g., TCGA-KIRC). CT enhancement phase labels, however, are inconsistent and often require manual inspection of DICOM metadata. Limited clinical metadata is available for some collections.

KiTS19, KiTS21, and KiTS23 focus on kidney cancer segmentation~\cite{kitsIntro}. KiTS23 provides increased diversity by including both corticomedullary and nephrogenic phase CT scans. All KiTS datasets supply high-quality 3D segmentations and basic clinical variables (e.g., age, sex, laterality, subtype), though not molecular data. Despite containing around $500$ cases, KiTS23 is considered large by medical imaging standards, where expert annotation is costly. Previous work has shown that state-of-the-art models such as nnU-Net, TransUNet, and SwinUNETR achieve strong generalisation using only a few hundred cases~\cite{isensee2021nnu}.

KiTS23 also offers rich multimodal information: voxel-level annotations, radiomics-feasible segmentations, clinical metadata, and extractable deep feature representations. Each CT volume can be decomposed into numerous overlapping 3D patches~\cite{wolf2023self}, effectively expanding the training set. Transfer learning and self-supervised pretraining further mitigate dataset size limitations, with studies demonstrating performance gains over training from scratch~\cite{tajbakhsh2016convolutional,kim2022transfer,wolf2023self,azizi2023robust}. Extensive augmentation strategies (rotations, flips, intensity perturbation) provide additional regularisation, and ImageNet-pretrained models have been shown to perform robustly even on datasets of this scale~\cite{maccin2025kidneynext}. Collectively, these factors support the suitability of KiTS23 for reliable DL classification research.

Additionally, portions of the TCGA imaging data within TCIA have been annotated as part of the AIAM initiative, where AI-generated segmentations were reviewed by expert radiologists~\cite{van2023image}. These standardised and quality-controlled annotations enhance the utility of TCIA for external validation. Consequently, TCIA datasets provide a strong foundation for testing the robustness and generalisability of models trained on KiTS23 across diverse real-world imaging conditions.

\subsection{Tumour Subtype Classification Performance}

Four clinically relevant tasks were identified across the included studies: (1) benign vs.\ malignant discrimination; (2) RCC subtype classification; (3) AML vs.\ RCC differentiation; and (4) oncocytoma vs.\ chromophobe RCC classification. A quantitative performance summary is provided in Table~\ref{tab:task_performance_final}. While moderate to high accuracy was reported in some settings, performance remains insufficient for routine clinical use, highlighting the potential for DL methods to improve robustness and subtype-specific discrimination.


\begin{table}[ht]
\centering
\caption{Summary of Model Performance by Classification Task}
\small
\begin{tabular}{lccccc}
\hline
\textsf{Classification Task} & \textsf{N Studies} & \textsf{Mean AUC} & \textsf{Std Dev} & \textsf{Min AUC} & \textsf{Max AUC} \\
\hline
Benign vs Malignant             & 10 & 0.809 & 0.076 & 0.715 & 0.960 \\
RCC Subtype Classification      & 6  & 0.814 & 0.068 & 0.720 & 0.910 \\
AML (wvf/fp) vs RCC             & 4  & 0.871 & 0.066 & 0.810 & 0.955 \\
Oncocytoma vs Chromophobe RCC  & 2  & 0.847 & 0.193 & 0.710 & 0.983 \\
\hline
\end{tabular}
\label{tab:task_performance_final}
\end{table}
\subsection{Systematic Review Conclusion}
Radiomics and ML approaches hold considerable promise for enhancing the non-invasive characterisation of renal tumours using CT imaging. This systematic review highlights a growing body of evidence supporting the feasibility of classifying tumour subtypes and distinguishing benign from malignant lesions with high diagnostic performance. Despite this progress, the field remains methodologically heterogeneous, with variation in classification tasks, feature selection techniques, imaging phases, and evaluation standards. Most studies employed traditional ML algorithms and handcrafted radiomics features, with limited adoption of DL. Corticomedullary and nephrographic phases were most commonly used, reflecting their diagnostic relevance and radiological contrast advantages.

Although several high-quality, large-scale public datasets such as KiTS, TCGA, and TCIA are available, featuring well-annotated, multi-phase CT scans ideal for radiomics research, relatively few studies have utilised them. This underuse limits the reproducibility and generalisability of current findings and slows the pace of clinical translation.

Future research should prioritise the use of larger, multi-institutional public datasets, standardised radiomics pipelines, and increased adoption of DL methods. With greater methodological rigour and data transparency, radiomics has the potential to become an integral component of renal tumour diagnosis and personalised management.

\section{CT-Based Radiomics Using PyRadiomics}

Radiomics enables the extraction of high-throughput quantitative features from medical images, converting them into structured variables for statistical or ML analysis. In renal tumour CT imaging, radiomics provides a non-invasive means of assessing tumour heterogeneity, intensity patterns, and morphology. Among existing radiomics tools, \textit{PyRadiomics} is widely used due to its standardisation, flexibility, and integration with Python~\cite{Griethuysen:2017}.

A typical workflow begins with acquisition of contrast-enhanced CT—most often in the corticomedullary or nephrographic phase—followed by pre-processing steps such as voxel resampling, intensity discretisation, and optional denoising. Segmentation then defines the region of interest, performed manually, semi-automatically, or using DL. PyRadiomics requires spatially aligned image–mask pairs in standard formats (e.g.\ NIfTI, DICOM).

From the segmented ROI, PyRadiomics extracts first-order statistics, shape descriptors, and texture features derived from matrices such as GLCM, GLRLM, GLSZM, NGTDM, and GLDM. Additional filtered features (e.g.\ wavelet, LoG) enable multi-scale texture analysis. Because this process can generate several hundred features per case, feature selection is essential. Common methods include variance and correlation filtering, RFE, LASSO, and mutual-information-based selection~\cite{VS, RFE}. Standardisation (e.g.\ z-scoring) is typically required before modelling.

Selected features are used to train ML classifiers such as SVM, random forests, logistic regression, or gradient boosting, with performance evaluated via accuracy, AUC, sensitivity, and specificity under cross-validation. PyRadiomics adheres to IBSI guidelines, facilitating reproducibility, and integrates seamlessly with Python ML libraries such as \texttt{scikit-learn}.

Despite these advantages, radiomics is sensitive to segmentation variability, acquisition differences, and feature instability. These limitations highlight the need for standardised imaging protocols, robust validation, and feature stability assessment.

\section{Computer Vision in Medical Imaging}
In cancer imaging, AI models are generally organ- or task-specific rather than universal, reflecting both tumour heterogeneity and the difficulty of modelling diverse anatomical contexts. Variations in histopathology, imaging appearance, and modality physics, together with domain shifts from scanners or acquisition protocols, hinder cross-domain deployment~\cite{bian2025artificial,ma2024segment,guan2021domain,koh2022artificial}. Public datasets are likewise organ-bound and rarely include cross-organ labels, constraining efforts toward universal detectors~\cite{wei2025multicenter,ali2021state}. Consequently, domain-specific systems remain standard, as they leverage anatomical priors and avoid the optimisation conflicts seen in multitask models~\cite{bian2025artificial,guan2021domain}. Clinical adoption also favours such models because their outputs align more closely with established diagnostic reasoning, whereas general-purpose systems often lack interpretability~\cite{ma2024segment,zhang2024challenges}.



\section{Fusion Strategies for Multimodal Data in Medical Imaging}

Multimodal DL integrates imaging, clinical, pathological, and genomic information to improve diagnostic and prognostic performance~\cite{huang2020fusion,mohsen2022artificial}. Reviews consistently show that combining modalities yields measurable gains: Kline \textit{et al.} reported a 6.4% accuracy increase with multimodal fusion~\cite{kline2022multimodal}, and Mohsen \textit{et al.} found early-fusion models outperform image-only or clinical-only baselines in 13 of 14 studies~\cite{mohsen2022artificial}. Similar benefits have been demonstrated across oncology, neurology, systemic disease, and treatment-response prediction~\cite{mohsen2022artificial,Vanguri2022,acosta2022multimodal}.

Recent surveys highlight rapid expansion of this field: Sun \textit{et al.} identified 77 multimodal DL studies using combinations of radiology, text, and clinical variables~\cite{sun2023scoping}, while Haq \textit{et al.} reported a sixteen-fold rise in multimodal radiology publications since 2019~\cite{haq2025mmml}. New reviews outline methodological trends—including integration with pathology and omics~\cite{Tariq2025,li2024review,simon2025multimodal}—and show that multimodal AI is moving toward clinically aligned systems~\cite{huang2020fusion,kline2022multimodal}. This thesis focuses on DL-based fusion strategies for recognition and classification tasks, covering early, intermediate, and late fusion~\cite{mohsen2022artificial,stahlschmidt2022}, along with VLM-based approaches and loss-function considerations~\cite{huang2020fusion,kline2022multimodal,li2024review}.

\subsection{Early, Intermediate and Late Fusion}

\textbf{Early fusion.} Early fusion concatenates raw or embedded features from multiple modalities into a unified representation~\cite{stahlschmidt2022,mohsen2022artificial}. Radiomics–clinical concatenation dominated early multimodal imaging research~\cite{mohsen2022artificial}, and several studies demonstrated clear improvements over unimodal models~\cite{huang2020fusion,kharazmi2018feature,yap2018multimodal}. Despite challenges such as modality imbalance~\cite{mohsen2022artificial}, early fusion remains widely used for tumour detection and risk prediction~\cite{yala2019deep,sun2025deep,xiao2025novel,hsu2021deep}.

\textbf{Intermediate fusion.} Intermediate (joint) fusion merges modality-specific subnetworks at a mid-level layer, enabling shared representation learning~\cite{mohsen2022artificial}. Examples include MRI–text fusion for tumour classification~\cite{he2021deep} and CT–clinical fusion for embolism severity prediction~\cite{grant2021deep}. Attention and probabilistic mechanisms are common~\cite{samak2020prediction,liu2022machine,kawahara2018seven}. Transformer-based multimodal models such as IRENE and cross-attention frameworks further enhance performance~\cite{khader2023multimodal,zhou2023transformer,khader2023medical}. GNN-based fusion has also emerged for integrating anatomical or genomic graphs with imaging~\cite{simon2025multimodal}. Intermediate fusion often outperforms both early fusion and unimodal models~\cite{mohsen2022artificial}.

\textbf{Late fusion.} Late fusion ensembles independently trained unimodal models~\cite{mohsen2022artificial}. Although it does not explicitly model cross-modal interactions, it is effective when modalities are asynchronous or incomplete. Demonstrated successes include Alzheimer’s diagnosis~\cite{qiu2018fusion}, outcome prediction with image+EHR ensembles~\cite{huang2020multimodal}, and stacked ultrasound–mammography–MRI classifiers~\cite{lu2024deep}. Hybrid late-fusion schemes can improve robustness under missing data~\cite{xiao2025novel}. It remains a transparent baseline for comparing integrated fusion strategies~\cite{mohsen2022artificial,huang2020multimodal}.

\subsection{VLMs and Advanced Multimodal Fusion Approaches} \label{sec:VLM}

Recent advances in multimodal learning have been driven by VLMs and related foundation models that jointly embed images and text via large-scale self-supervised pretraining~\cite{ryu2025vlmreview}. Inspired by CLIP~\cite{radford2021learning}, medical VLMs such as BioViL~\cite{boecking2022biovil,lyu2024language}, GLoRIA~\cite{huang2021gloria}, MedCLIP~\cite{park2020contrastive,wang2022medclip}, and BioMedCLIP~\cite{zhang2023biomedclip} align chest X-rays or other medical images with associated reports, enabling zero-shot classification and image–text retrieval. Earlier multimodal networks like TieNet had already demonstrated that incorporating clinical text improved chest X-ray classification by 5–6\% AUC~\cite{wang2018tienet}. Surveys now catalogue over 60 medical foundation models across radiology, pathology, and ophthalmology~\cite{ryu2025vlmreview}, including encoder–decoder systems for report generation and visual question answering~\cite{hartsock2024vlmreport,thawkar2023xraygpt}.

A growing line of work “textualises” structured or numeric data for VLM input. AutoRad-Lung converts radiomics features into descriptive prompts for lung CT malignancy prediction~\cite{khademi2025autorad}, and FM-Bridge encodes imaging and risk variables as natural-language statements to improve liver tumour diagnosis~\cite{huang2025fmbridge}. Prompt-based methods similarly adapt general VLMs to domain-specific tasks, as in dental X-ray analysis~\cite{du2024prompting}. In pathology, MIL frameworks such as MI-Zero and ViLA-MIL use CLIP-based embeddings to align gigapixel slide patches with histological text, enabling zero-shot cancer subtype classification~\cite{lu2023mizero,shi2024vila}.

Key challenges for medical VLMs include limited domain data, specialised terminology, and factual reliability~\cite{ryu2025vlmreview}. Proposed solutions span domain vocabularies, better region–word alignment, and expert feedback. MedVLM-R1 uses RLHF to refine reasoning~\cite{pan2025medvlm}, and SERPENT-VLM incorporates factuality checking to reduce hallucinations~\cite{kapadnis2024serpent}. Large annotated datasets such as MIMIC-CXR and ROCO support ongoing training and transfer~\cite{huang2020fusion,haq2025mmml}. Emerging MLLMs extend these capabilities to interactive, image–text understanding, and compact variants (e.g.\ Small-MM) show that clinically deployable models need not be excessively large~\cite{nam2025multimodal,zambrano2025smallmm}. Overall, VLM-based approaches point toward holistic, patient-centred models that integrate visual, textual, and molecular information~\cite{haq2025mmml,zhoudeep}.

Loss design in multimodal DL balances task performance with cross-modal alignment. Most models still rely on standard supervised objectives—cross-entropy for classification, Dice or pixel-wise cross-entropy for segmentation, and MSE for regression~\cite{mohsen2022artificial}. Class imbalance is commonly addressed with weighting or focal loss~\cite{kline2022multimodal,haq2025mmml}, multi-label problems with sigmoid cross-entropy, and multi-task settings with weighted combinations of classification and regression terms~\cite{huang2020fusion}. Auxiliary modalities usually enhance representations rather than change the primary loss, although reconstruction or imputation losses can regularise models when modalities are missing~\cite{haq2025mmml}.

Self-supervised and contrastive frameworks introduce additional objectives to enforce cross-modal consistency. InfoNCE and CLIP-style losses maximise similarity between paired image–text samples while minimising mismatched pairs~\cite{windsor2023vision,ryu2025vlmreview}. Medical adaptations such as BioViL and GLoRIA incorporate both global and local alignment terms, including local ranking or correlation-based losses for aligning image regions with textual phrases~\cite{huang2021gloria}. Reconstruction losses can further encourage one modality to predict another, regularised by $L_2$ penalties~\cite{stahlschmidt2022,krones2025review}. In survival analysis, Cox partial likelihood is used to model time-to-event risk in multimodal settings~\cite{huang2020fusion}, and hybrid tasks (e.g.\ segmentation plus classification) employ composite losses to balance objectives~\cite{oltu2025automated,hasan2024recent}. Consistency- and distillation-based penalties have also been proposed for handling missing modalities and improving robustness~\cite{huang2020fusion,fang2024aligning}.

Recent models integrate more advanced objectives, including RL for multimodal generation (e.g.\ MedVLM-R1)~\cite{pan2025medvlm} and joint alignment-plus-cross-entropy schemes to reduce hallucinations in VLM outputs~\cite{ryu2025vlmreview}. Adversarial and distillation losses have been explored for realistic report generation and missing-modality learning~\cite{fang2024aligning}. Overall, the main challenge lies less in inventing new losses than in appropriately weighting existing ones and ensuring effective gradient flow through all modality branches~\cite{mohsen2022artificial}.

\subsection{Summary} \label{sec:conclusion-lit}

In summary, multimodal fusion strategies in medical imaging span early, intermediate, and late integration, each offering distinct trade-offs in complexity, performance, and interpretability~\cite{mohsen2022artificial,kline2022multimodal}. Early fusion remains simple and widely adopted, often outperforming unimodal baselines~\cite{mohsen2022artificial,kline2022multimodal}. Intermediate fusion, particularly attention- and transformer-based designs, enables joint feature learning and frequently delivers state-of-the-art results~\cite{mohsen2022artificial,zhou2023transformer,simon2025multimodal}. Late fusion serves as a practical, interpretable benchmark and can be advantageous when modalities are heterogeneous or asynchronously available~\cite{mohsen2022artificial,huang2020multimodal,lu2024deep}.

A major recent development is the emergence of VLMs and multimodal foundation models, which align images, text, and textualised structured data within shared latent spaces~\cite{ryu2025vlmreview,radford2021learning,boecking2022biovil,huang2021gloria,park2020contrastive,zhang2023biomedclip}. Approaches such as AutoRad-Lung and FM-Bridge suggest that converting radiomics and clinical variables into descriptive text may be a powerful strategy for leveraging VLMs in medical imaging~\cite{khademi2025autorad,huang2025fmbridge}. This thesis builds on these ideas by applying textualisation and VLM-based fusion to 3D renal CT, treating the VLM as a multimodal feature extractor for supervised subtype classification. Loss-function design in multimodal learning typically combines standard supervised objectives with contrastive or consistency terms to ensure effective cross-modal alignment~\cite{windsor2023vision,huang2021gloria,haq2025mmml}. Across tasks and modalities, meta-analyses show that multimodal models reliably outperform unimodal counterparts when properly optimised~\cite{kline2022multimodal,haq2025mmml,simon2025multimodal,huang2020fusion,mohsen2022artificial}.

Persistent challenges include missing modalities, robustness and fairness across populations, and the need for interpretable models~\cite{kline2022multimodal,YANG202229,haq2025mmml}. Explainable AI tools such as attention visualisation and prototype-based reasoning can clarify modality contributions and improve clinician trust~\cite{YANG202229,diao2025ppal}. Future directions focus on unified architectures for images, text, and structured data; integration of large language models for medical reasoning; and efficient, distilled multimodal systems suitable for deployment~\cite{zambrano2025smallmm,Vanguri2022,zhoudeep}. More broadly, multimodal DL is evolving towards holistic, interpretable, and clinically grounded AI systems that better reflect the complexity of real-world patient care.

\section{Interpretability and Explainability in Medical AI}

In safety-critical settings such as medical imaging, interpretability—the ability to understand how a model reaches its predictions—is increasingly seen as essential for clinical adoption~\cite{abgrall2024should,chen2022explainable}. Surveys indicate that many clinicians remain sceptical of “black-box” AI systems; for example, 71\% of ICU clinicians report uncertainty about AI reliability~\cite{abgrall2024should}. Interpretability helps bridge this trust gap by allowing experts to verify that AI reasoning aligns with medical knowledge and ethical standards, and it supports emerging regulatory requirements such as the EU AI Act~\cite{abgrall2024should}. While often conflated, interpretability differs from post hoc explainability: interpretable models are inherently transparent (e.g.\ linear models, decision trees)~\cite{houssein2025explainable}, whereas explainability methods (e.g.\ LIME, SHAP) approximate the logic of complex models without fully revealing their internal workings~\cite{houssein2025explainable}. Recent accountable-AI frameworks advocate prioritising inherently interpretable components and using post hoc explanations as complementary tools~\cite{singh2025beyond,napravnik2025lessons}.

Empirical studies show that well-designed explanations improve trust calibration and decision quality~\cite{chen2022explainable,vani2025personalized}. Heatmaps and saliency maps help radiologists assess whether a model attends to clinically relevant regions~\cite{vani2025personalized,napravnik2025lessons}, and pretraining on radiology-scale data can yield more meaningful visual explanations~\cite{napravnik2025lessons}. Explanations must remain human-centred—understandable, contextual, and clinically actionable~\cite{chen2022explainable}—especially in oncology, where opaque reasoning remains a major barrier despite high reported accuracies~\cite{hrinivich2023interpretable}. Interpretability also supports bias detection and error analysis before patient harm occurs~\cite{abgrall2024should}. Accordingly, recent work argues for explicitly trading a small amount of accuracy for greater transparency when necessary~\cite{singh2025beyond}.

For multimodal medical AI, interpretability is increasingly viewed as a practical necessity rather than an optional add-on~\cite{napravnik2025lessons,vani2025personalized}. While post hoc explainability remains useful for complex models, recent work emphasises designing transparency into model architectures from the outset~\cite{singh2025beyond}. This includes choosing structures amenable to explanation (e.g.\ attention mechanisms, case-based reasoning, generative rationales) and ensuring outputs are human-verifiable~\cite{sharma2025xai}. Inherently interpretable models, which combine explicit risk factors with imaging evidence in transparent ways, tend to align better with clinical reasoning and support more reliable decision-making than opaque black-box systems~\cite{houssein2025explainable,vani2025personalized,jungmann2023algorithmic}. Interpretability also involves communication: dashboards, interactive visualisations, and query-based interfaces can surface AI reasoning in clinician-friendly formats~\cite{pahud2024orchestrating,sharma2025xai}. Achieving an appropriate balance between performance and interpretability will be crucial as multimodal AI transitions from research into routine clinical use~\cite{vani2025personalized,chen2022explainable}.

\subsection{Interpretability across fusion strategies}

Explaining multimodal models is particularly challenging because they must reveal not only which features influence predictions but also how modalities interact. Early fusion tends to entangle features from all inputs, making it difficult to attribute importance to specific modalities. In such models, explanations rely on global feature-importance or perturbation-based methods (e.g.\ SHAP)~\cite{sharma2025xai,vani2025personalized}, as well as occlusion experiments to quantify each modality’s impact~\cite{vani2025personalized}. However, high-dimensional image embeddings may dominate smaller feature sets, reducing clarity and leaving these models effectively black boxes.

Intermediate fusion can improve interpretability by keeping modality-specific branches separate until mid-level integration~\cite{stahlschmidt2022}. This structure allows the use of modality-specific XAI techniques—e.g.\ Grad-CAM for imaging and SHAP for clinical features—followed by inspection of fusion weights or attention maps to assess cross-modal contributions~\cite{sharma2025xai}. Attention-based fusion may align image regions with text tokens, more closely reflecting radiological reasoning~\cite{sharma2025xai}. Gating mechanisms that assign instance-specific modality weights offer additional transparency~\cite{sharma2025xai}. Nonetheless, attention weights do not automatically guarantee faithful explanations and require validation~\cite{napravnik2025lessons,pahud2024orchestrating}.

Late fusion, which combines separate unimodal predictions, offers the greatest modality-wise transparency but limited insight into interactions. Clinicians can interpret each branch independently, using heatmaps for imaging and feature charts for clinical data, mirroring multidisciplinary reasoning~\cite{abgrall2024should}. Ensemble-level XAI methods can approximate cross-modal effects~\cite{sharma2025xai}. While late fusion may sacrifice some representational synergy, its modularity and clarity make it attractive in safety-critical workflows~\cite{abgrall2024should}.

\subsection{Vision–language and foundation models}

VLMs naturally support more intuitive explanations by operating directly on images and text. Rather than exposing only abstract feature vectors, they can generate natural-language justifications for predictions, aligning AI reasoning with radiological narratives~\cite{sharma2025xai,pahud2024orchestrating}. Attention maps in CLIP-style models reveal associations between image regions and semantic tokens, while architectures such as GLoRIA explicitly link patches to descriptive words~\cite{napravnik2025lessons,huang2021gloria}. Textualisation of non-image inputs (e.g.\ lab values, genomics, radiomics) further enhances interpretability by mapping heterogeneous data into a shared language space~\cite{sharma2025xai}. However, VLMs can hallucinate or overgeneralise, motivating the use of factuality constraints and RLHF to ensure clinically reliable explanations~\cite{pahud2024orchestrating}. As multimodal foundation models mature, they hold promise for systems that combine strong predictive performance with human-readable, clinically meaningful explanations.

\subsection{Clinician acceptance and evaluation of XAI}

Clinician acceptance is central to effective interpretability. Studies indicate that personalised, case-specific explanations improve trust and support appropriate reliance on AI recommendations~\cite{vani2025personalized}. Vani \textit{et al.} showed that SHAP- and attention-based visualisations enhanced physicians’ understanding of model outputs~\cite{vani2025personalized}, while other work argues that explainable AI can increase guideline adherence when explanations are clear and contextually relevant~\cite{abgrall2024should,houssein2025explainable}. Poor or misleading explanations, by contrast, may reduce safety and confidence~\cite{abgrall2024should,napravnik2025lessons}. Consequently, a “clinician-in-the-loop” approach is increasingly advocated, where experts evaluate AI explanations through user studies comparing diagnostic performance with and without explanations~\cite{napravnik2025lessons,vani2025personalized}. Standards for clinical explainability, such as understandability, truthfulness, and actionability, are being proposed~\cite{pahud2024orchestrating,jin2023guidelines}, and professional societies and regulators now highlight interpretability as a key requirement for future AI approval~\cite{abgrall2024should,houssein2025explainable}.





\section{Research Gaps}

First, renal CT imaging research still relies predominantly on handcrafted radiomics with conventional ML classifiers, with limited use of DL and few controlled comparisons between the two paradigms. Reported performances vary widely owing to differences in imaging phase, feature extraction, and evaluation criteria, leading to poor reproducibility and generalisability.

Second, although high-quality public datasets such as KiTS23 and TCGA/TCIA exist, they are not yet systematically exploited, and external validation remains uncommon. Standardised pipelines and benchmark evaluations on public data are needed to enable robust comparison across studies.

Third, while prior work highlights the potential of early, intermediate, and late fusion, as well as transformer-based and VLM approaches, renal tumour imaging has seen few systematic comparisons of these multimodal strategies or analyses of how fusion mechanisms affect accuracy, robustness, and interpretability. In particular, textualising radiomics features and aligning them with images within a VLM framework remains largely unexplored, despite its promise for improving clinical transparency and human-understandable reasoning.

Fourth, most existing studies omit patient-level metadata that could enrich the multimodal context, and rarely quantify robustness to missing modalities or interpretability aligned with radiologist reasoning. Without explicit assessment of how models communicate their decision rationale, the translational value and clinical trustworthiness of multimodal AI remain uncertain.

Finally, although recent work has proposed various multimodal loss functions, task-specific strategies for balancing supervised objectives with cross-modal alignment and consistency losses in 3D renal CT classification are still underexplored.

In summary, this thesis addresses these gaps by: (1) establishing reproducible baselines with external validation on public 3D CT datasets; (2) systematically comparing multiple fusion strategies, including transformer-based and textualised VLM-based integration of radiomics; (3) incorporating patient metadata to build tri-modal models; and (4) analysing loss-function designs that balance supervised and cross-modal learning. Together, these directions aim to improve fusion efficacy, robustness, and interpretability, and to support clinically trustworthy, explainable multimodal AI for renal tumour analysis.

Based on the research gaps explored in the literature review, this thesis aims to answer the  research questions mentioned in section~\ref{rq}.
